\documentclass[11pt,a4paper]{article}

% ------------------------------------------------------------------
%  v3: repositioned as a standalone proposal paper.
%  Compiles with pdfLaTeX (TeXLive / MiKTeX) or XeLaTeX / Tectonic.
%  Target journals (in order of realism):
%    1. Foundations of Physics  (Springer)
%    2. Classical and Quantum Gravity / IJMPD
%  For submission, journals provide their own class files; the content
%  below ports directly (no exotic macros used).
% ------------------------------------------------------------------

\usepackage[T1]{fontenc}
\usepackage[utf8]{inputenc}
\usepackage{lmodern}
\usepackage{amsmath,amssymb}
\usepackage{graphicx}
\usepackage{booktabs}
\usepackage[margin=2.5cm]{geometry}
\usepackage{microtype}
\usepackage[numbers,sort&compress]{natbib}
\usepackage[colorlinks=true,linkcolor=blue,citecolor=blue,urlcolor=blue]{hyperref}

\newcommand{\meq}{MEQ}
\newcommand{\um}{\,\mu\mathrm{m}}

\title{The Minimum Energy Quantum (\meq{}) Framework:\\
A Single-Quantum Kinetic Substrate for the Unification of\\
Quantum Field Theory and Gravitation, with Four Falsifiable Predictions}

\author{Roger Wang (Wang Yanxiang)\\[2pt]
\small Tsinghua x-lab, Shenzhen, China; ALAYA Tech, Beijing, China\\
\small Correspondence: wanyanx14@tsinghua.org.cn}
\date{\today}

\begin{document}
\maketitle

\begin{abstract}
\noindent
We propose the Minimum Energy Quantum (\meq{}) framework: a unified account in which all observable
matter and interactions emerge as collective excitations of a single species of massless quanta
moving at the speed of light, interacting through elastic exchange. The framework replaces the
independent postulate sets of quantum field theory (QFT) and general relativity (GR) with one
microscopic substrate, and contributes four novel elements: (i) an operational derivation of the
relativity postulates from the light-speed substrate --- the invariance of $c$ is a fixed point of
any measurement performed ``with $c$ on $c$''; (ii) an induced-gravity self-consistency structure,
in which the Einstein equation arises as the leading term of a generally covariant mean-field
condition on the substrate, with higher-curvature corrections suppressed by a substrate length
$\ell_\ast$; (iii) topological emergence of fermionic quasiparticles from bosonic constituents via
the Skyrme/Goldstone--Wilczek mechanism; and (iv) emergent --- not exact --- Lorentz symmetry. We
derive the complete set of structural requirements and quantitative bounds that define the viable
parameter space: exact elasticity with detailed balance (no heating), Lorentz-invariant induced
interactions (orbital stability; $c_g = c$ from GW170817), universal coupling through
$T_{\mu\nu}$ (validated by MICROSCOPE at $10^{-15}$), a spin-2 collective gravitational quantum
(consistent with post-Newtonian phenomenology and tensor gravitational waves), and cosmological
transparency ($\xi \lesssim 10^{-3}$ fractional photon energy loss per Hubble length). Within this
viable region the framework makes four falsifiable predictions, the first three supplied with full
quantitative estimates and discovery reach: (1) a Yukawa-type deviation from the inverse-square law
for $\ell_\ast \lesssim 40\um$, with the reach of next-generation torsion-balance and
micro-cantilever experiments mapped in the $(\alpha,\lambda)$ plane (Sec.~\ref{sec:p1} and
Appendix~\ref{app:quant}); (2) an anomalous vacuum-birefringence coefficient $\zeta$ below the
Euler--Heisenberg floor, with the reach of optical-cavity experiments computed for PVLAS-class and
next-generation strong-field setups; (3) a mass non-stationarity of gravitating configurations ---
a slow drift of inertial mass scaling as $\dot m/m = \eta\, g/c$ --- for which we derive existing
bounds from binary-pulsar orbital decay ($\eta \lesssim 5\times10^{-8}$, improving to
$\lesssim 9\times10^{-10}$ for the double pulsar), mass metrology, and atomic-clock comparisons,
and we specify a differential experiment with annual modulation tagging capable of probing
$\eta_{\mathrm{diff}} \sim 10^{-8}$--$10^{-11}$; and (4) phase-dependent galactic gravitation,
stated as a target for high-redshift surveys. The framework is thereby exposed to decisive
evidence within the current decade.
\end{abstract}

\noindent\textbf{Keywords:} quantum gravity; emergent gravity; induced gravity; kinetic theory;
equivalence principle; dark matter; experimental tests

\tableofcontents

\section{Introduction}\label{sec:intro}

\subsection{The problem}

Quantum field theory and general relativity are the two experimentally validated pillars of
fundamental physics, yet they rest on mutually incompatible foundations: QFT presumes a fixed
background spacetime and local quantum states defined on it, while GR makes the metric itself
dynamical. Direct perturbative quantization of the Einstein--Hilbert action fails (the theory is
power-counting non-renormalizable beyond one loop), while the major structured programs ---
string/M-theory, loop quantum gravity, asymptotic safety, causal dynamical triangulations,
holography --- have not produced a uniquely selected, experimentally verified low-energy
phenomenology \cite{Rovelli2015,Polchinski1998,Reuter1998,Loll2020,Batista2025}.

Three empirical facts sharpen the problem rather than resolving it:
\begin{enumerate}
\item \textbf{Gravity couples to all energy.} The MICROSCOPE mission has verified the
universality of free fall to $|\eta| = (-0.53 \pm 2.05)\times10^{-15}$ \cite{Touboul2022}, and
lunar laser ranging confirms that gravitational binding energy itself gravitates
\cite{Williams2004} --- gravity is nonlinear and self-coupling.
\item \textbf{The graviton, if it exists, is light-speed and effectively spin-2.} GW170817
constrained $|c_g - c|/c \lesssim 10^{-15}$ \cite{GW170817}, and the orbital dynamics measured by
LIGO/Virgo agree with the tensor (spin-2) polarization structure of GR, disfavoring large scalar
or vector gravitational radiation components.
\item \textbf{Photon propagation is extraordinarily clean.} No energy loss, no blurring, and no
dispersion attributable to the propagation medium has ever been detected, down to fractional
sensitivities approaching the Planck scale in high-energy gamma-ray timing
\cite{Abdo2009,Vasileiou2013,Myers2005}.
\end{enumerate}
These three facts are precisely the boundary conditions that any unified substrate must satisfy.
The proposal of this paper is built to satisfy them \emph{by construction}, and to convert them
into quantitative design constraints that locate the theory's novel, testable content.

\subsection{Historical lineage}

The idea that gravity arises from the collective behavior of a more fundamental medium has a long
and productive history:
\begin{itemize}
\item \textbf{Le Sage (1784)} and Fatio de Duillier before him proposed ``push gravity'' from
an isotropic flux of ultramundane corpuscles, with bodies shielding each other from the flux.
Maxwell and Poincar\'e showed that \emph{inelastic} interception implies catastrophic heating (for
Earth, $\geq 10^{59}\,$W), while Laplace-type aberration arguments show that a \emph{non-covariant}
retarded flux destabilizes orbits \cite{LeSage2002,Poincare1908}. Both objections are statements
about \emph{which} microscopic dynamics are admissible, and both are evaded by the design choices
of Sec.~\ref{sec:kinetic}.
\item \textbf{Sakharov (1968)} showed that the Einstein--Hilbert action can be interpreted as the
elastic response of the quantum vacuum --- gravity as an \emph{induced} phenomenon of the matter
vacuum energy \cite{Sakharov1968}. This survives as the effective-field-theory (EFT) view of
quantum gravity \cite{Donoghue1994}.
\item \textbf{Jacobson (1995)} derived the Einstein equation as an equation of state of local
Rindler horizons \cite{Jacobson1995}; \textbf{Verlinde (2010)} recast gravity as entropic
\cite{Verlinde2011}; \textbf{Volovik} and the analog-gravity program demonstrated that Lorentzian
spacetime, horizons, and effective metrics arise generically as low-energy phenomena of
condensed substrates (superfluid $^3$He) \cite{Volovik2003,Barcelo2005}.
\item \textbf{Superfluid vacuum theories} (Sinha, Sivaram, Sudarshan; Zloshchastiev) treat
particles as quasiparticles of a superfluid vacuum \cite{Sinha1976,Zloshchastiev2011}.
\end{itemize}
The \meq{} framework proposed here belongs to this lineage and is distinguished by four
commitments:
\begin{itemize}
\item \textbf{(P1) Single species.} Only one type of fundamental quantum exists, the ``minimum
energy quantum'' (\meq{}), all moving at $c$ and interacting only through elastic collisions.
\item \textbf{(P2) Matter as bound states.} Bosons are coherent groups of co-moving \meq{}s
(vector-field excitations); fermions are topological/rotational bound configurations (spinor
fields) of the substrate.
\item \textbf{(P3) Gravity as collective field.} The gravitational field is not produced
\emph{by} matter but arises \emph{with} matter, as the self-consistent collective field of the
substrate surrounding every bound configuration.
\item \textbf{(P4) Emergent spacetime.} Time and space, as measured quantities, inherit their
properties (minimum scales, Lorentz symmetry, curvature near gravity) from the light-speed
substrate.
\end{itemize}

\subsection{What this paper does}

The purpose of this paper is threefold:
\begin{enumerate}
\item \textbf{Formulate} the framework (P1)--(P4) as a mathematically explicit relativistic
kinetic mean-field theory with an induced effective action (Sec.~\ref{sec:kinetic}), making
every assumption explicit.
\item \textbf{Establish consistency} with the complete set of modern experimental constraints,
deriving quantitative bounds on the \meq{}--matter and \meq{}--photon interaction cross-sections
that define the viable parameter space (Sec.~\ref{sec:exp}).
\item \textbf{Extract the falsifiable content}: a parameterized family of deviations from
GR/Newtonian gravity and QED, with target experiments identified (Sec.~\ref{sec:pred}), and full
quantitative estimates for the first three predictions (Appendix~\ref{app:quant}).
\end{enumerate}

This is a hypothesis paper with an explicit research program; we do not claim a completed
derivation of the Standard Model from the substrate. We claim that the \meq{} framework is
\emph{consistent with all existing data} in a well-defined region of parameter space, and that it
makes \emph{distinctive, testable} predictions in that region --- predictions which no
established framework (GR, QED, or their EFT extensions) currently makes in the same form.

\section{Conceptual Foundations}\label{sec:concept}

\subsection{Operational spacetime}

A central innovation of the framework is its operationalist account of spacetime: all
measurements of space and time are performed by apparatus ultimately composed of light-speed
quanta. If every clock is ultimately a light-clock and every ruler is delimited by signals
propagating at $c$, then:
\begin{enumerate}
\item \textbf{Light-speed invariance is not an additional postulate} but a fixed point of the
measurement process itself --- ``measuring $c$ with $c$'' returns $c$ in every frame. This
gives a concrete substrate realization of the observation underlying Mansouri--Sexl-type test
theories \cite{Mansouri1977} and doubly-special-relativity programs \cite{Amelino2002}: the
invariance of the limiting speed is structurally guaranteed when the substrate that defines
measurement is the same substrate that defines light.
\item \textbf{Minimum measurable scales} arise because the substrate quanta are not points. If
the \meq{} has an effective size $\ell_\ast$, then no apparatus built from \meq{}s can resolve
intervals smaller than $\mathcal{O}(\ell_\ast)$. The discrete spectra of geometric operators in
loop quantum gravity \cite{Rovelli2015} can be reinterpreted in this language: the minimum
area/volume is set by the substrate granularity.
\item \textbf{Curvature near gravity} is an operational statement: in a region where light
trajectories are bent, all measurements performed with light are distorted exactly as if the
metric were curved. The statement ``gravity \emph{is} spacetime curvature'' is derived in this
framework as ``curvature is how a light-speed substrate's collective field manifests in
light-based measurements.'' We regard (following \cite{Volovik2003}) the metric as an
\emph{effective} object; this is also the EFT position \cite{Donoghue1994}.
\end{enumerate}

\subsection{Postulates}

We state the \meq{} postulates in a form suitable for quantitative analysis:
\begin{itemize}
\item \textbf{A1 (Substrate).} There exists a universal quantum field $\psi$ whose quanta are
massless ($E = |\mathbf{p}|c$), with an interaction that reduces to elastic scattering at leading
order.
\item \textbf{A2 (Emergence of matter).} The observed particle spectrum (spin-0, $\tfrac12$, 1)
arises as collective excitations --- bound and/or topological configurations --- of $\psi$.
\emph{The task of Sec.~\ref{sec:fermions} is to state the known mechanisms by which bosonic
constituents yield fermionic excitations.}
\item \textbf{A3 (Emergence of gravity).} The gravitational field $g_{\mu\nu}$ is the collective
mean field of the substrate; it is determined by a self-consistency condition on the total
stress-energy, of which the Einstein equation is the leading-order term.
\item \textbf{A4 (EFT matching).} At energies well below the substrate scale
$E_\ast \sim \hbar c/\ell_\ast$, the framework must reduce to (and is constrained by) the
Standard Model plus GR, with corrections suppressed by powers of $E/E_\ast$.
\end{itemize}
A2 is the least developed link in the chain. We note two established mechanisms showing it is
not \emph{a priori} impossible: (i) in the Skyrme model, fermions emerge as topological solitons
of bosonic fields, with the fermionic statistics derived from the
Wess--Zumino--Goldstone term \cite{Skyrme1962,Goldstone1981}; (ii) in 2+1 dimensions,
flux--charge composites (anyons) interpolate between bosonic and fermionic statistics
\cite{Wilczek1982}. A1--A4 therefore demand a 3+1-dimensional analogue of such mechanisms; we
return to this in Sec.~\ref{sec:open}.

\subsection{Classification of observables}

We introduce a trichotomy of observables: \textbf{(i) directly detected quanta} (detected by
absorption in apparatus --- e.g., photons); \textbf{(ii) indirectly inferred quanta} (inferred
from scattering of substrate quanta off configurations --- ordinary fermionic matter);
\textbf{(iii) inferred fields} (gravitational, and historically, electromagnetic fields; and
hypothetical entities such as dark matter). The epistemic status differs across the three
classes, and any theory that promotes a class-(iii) entity to class-(i)/(ii) status --- or
eliminates it --- gains in economy. This is the sense in which the \meq{} program is, at its
core, a program in \textbf{theory economy}: replace the standard model's particle zoo plus the
graviton plus (possibly) dark matter with one substrate quantum.

\section{The Kinetic Framework}\label{sec:kinetic}

\subsection{Distribution function and conservation}

Let $f(x,\mathbf{n})$ denote the phase-space density of substrate quanta at spacetime point $x$
propagating in direction $\mathbf{n}$ (all quanta are on-shell at speed $c$). Its stress-energy
tensor is
\begin{equation}
T^{\mu\nu}(x) \;=\; \sigma\, c \int d\Omega\, n^\mu n^\nu\, f(x,\mathbf{n}),
\end{equation}
(with $\sigma$ an energy scale factor), and the dynamics is governed, at the microscopic level, by
a Boltzmann-type equation
\begin{equation}
n^\mu \partial_\mu f \;=\; C[f,F],
\end{equation}
where $C$ is the elastic collision functional and $F$ denotes the self-consistent mean field.
Elasticity implies exact local conservation of $T^{\mu\nu}$ after integration over directions ---
the substrate exchanges momentum and energy with matter configurations but never dissipates into
heat. This is the first structural difference from Le Sage-type theories: \textbf{the heating
catastrophe is absent because collisions are exactly elastic}, and the drag catastrophe is absent
in the limit where the anisotropic part of $f$ is of order the gravitational field itself.
We quantify both statements in Secs.~\ref{sec:heat}--\ref{sec:aberration}.

\subsection{Mean-field self-consistency: induced gravity}\label{sec:induced}

The central equation of the framework is a self-consistency condition between the substrate
stress-energy and the effective metric that the substrate's collective field defines for all
propagation (including its own quanta). At leading order in derivatives, the most general
generally covariant self-consistency condition is
\begin{equation}
G_{\mu\nu}[g] + \Lambda g_{\mu\nu} + \mathcal{O}(\ell_\ast^{2} R^{2})
\;=\; \frac{8\pi G}{c^{4}}\, \Big( \langle T_{\mu\nu}\rangle_{\rm substrate} + T_{\mu\nu}^{\rm matter} \Big),
\label{eq:einstein}
\end{equation}
which is the Einstein equation with higher-curvature corrections --- Sakharov's induced gravity
in modern EFT form \cite{Sakharov1968,Donoghue1994}. Within the \meq{} framework this equation
acquires a specific interpretation: \textbf{the left-hand side is the collective, continuum
response of the substrate; the right-hand side contains both ordinary matter (quasiparticle
excitations) and the substrate's own anisotropic background}. Two consequences follow:
\begin{enumerate}
\item \textbf{Nonlinearity is automatic.} Because the substrate's stress-energy includes its own
gravitational binding configuration, ``gravity gravitates'' --- the requirement imposed by lunar
laser ranging and the Nordtvedt effect \cite{Williams2004}. A purely linear superposition field
cannot satisfy this; the self-consistency structure is essential to the framework.
\item \textbf{Corrections are predicted to exist} at the scale $(\ell_\ast/\ell)^{2}$ for
curvature radius $\ell$. If $\ell_\ast$ is of order the Planck length, these are unobservably
small; if $\ell_\ast$ is larger (e.g., of order $10^{-5}$--$10^{-3}\,$m, near the reach of
short-range gravity tests), they are testable. The framework thus has a natural scanning
parameter, and the experimental situation (Sec.~\ref{sec:shortrange}) bounds it from below ---
this is the origin of Prediction 1 (Sec.~\ref{sec:p1}).
\end{enumerate}

\subsection{Fermionic excitations from bosonic constituents}\label{sec:fermions}

In our framework, a fermion is realized as a \textbf{topological soliton or flux--charge
composite} of the substrate, with fermionic statistics arising from a configuration-space
(braid) phase --- the Skyrme/Goldstone--Wilczek mechanism \cite{Skyrme1962,Goldstone1981}, which
\emph{derives} an odd fermion number and half-integer spin from bosonic fields with a topological
term. The \meq{} ``spinor field'' of the postulate A2 should therefore be read as such a
configuration; its quantitative construction --- a 3+1-dimensional realization covering the
observed chiral spectrum --- is the central open problem of the program (Sec.~\ref{sec:open}),
and no claim of completion is made here.

\subsection{Lorentz invariance from the substrate}

A massless substrate at equilibrium supports emergent Lorentz invariance for its quasiparticles,
in the sense of Volovik's classification \cite{Volovik2003}: the low-energy quasiparticle
dispersion $E^{2} = c^{2}p^{2}\bigl(1 + \mathcal{O}(p^{2}/E_\ast^{2})\bigr)$ with emergent (not
fundamental) Lorentz symmetry. The substrate picture \emph{predicts tiny violations} of exact
Lorentz invariance at energies of order $E_\ast$, which is testable:
\begin{itemize}
\item Gamma-ray burst photon timing constrains linear-in-$(E/E_\ast)$ dispersion to
$E_\ast \gtrsim E_{\rm Planck}$ for most operator structures \cite{Abdo2009,Vasileiou2013};
\item the absence of vacuum Cherenkov radiation and photon splitting for high-energy photons
bounds certain higher-order structures \cite{Jacobson2003}.
\end{itemize}
These are among the sharpest existing constraints on the \meq{} substrate scale, and they force
the substrate to be \emph{extraordinarily} --- but not exactly --- transparent and
dispersion-free. Parameterizing the \meq{}--photon forward-scattering amplitude in the standard
dimension-5/6 EFT notation \cite{Myers2005,Bolokhov2007} gives the explicit bounds quoted in
Sec.~\ref{sec:p2}.

\section{Experimental Consistency and Parameter-Space Constraints}\label{sec:exp}

Any unified substrate theory must satisfy a complete set of modern experimental constraints. In
this section we derive the structural requirements and quantitative bounds that define the viable
parameter space of the \meq{} framework. Each subsection states a requirement of the model, the
data that fixes it, and its quantitative consequence. Far from weakening the framework, these
constraints are what give it empirical content: they locate precisely where its novel predictions
can be sought (Sec.~\ref{sec:pred}).

\subsection{Thermodynamics: exact elasticity and detailed balance}\label{sec:heat}

The Maxwell--Poincar\'e heating objection \cite{LeSage2002,Poincare1908} --- that bodies should be
incinerated by a flux medium --- is evaded \emph{by construction} in the \meq{} framework:
collisions are exactly elastic (A1), so no energy is dissipated as heat. The dynamical mechanism
of gravity is then absorption \emph{into} bound configurations balanced by emission in detailed
balance: \meq{}s absorbed into a bound configuration \emph{become part of its mass}. The energy
bookkeeping is not ``heat'' but ``mass'': equilibrium requires that the rate of \meq{} absorption
equal the rate of mass-energy growth of the configuration, and that configurations in
equilibrium absorb and emit in detailed balance. \textbf{The framework therefore predicts that
isolated bodies are, at some level, not exactly stationary in mass} --- a falsifiable (and
currently untested at the required precision) feature, developed quantitatively as Prediction 3
(Sec.~\ref{sec:p3}). A strict equilibrium limit exists in which absorption and emission balance
exactly, and gravity is the residual anisotropic stress of this balanced exchange --- the kinetic
counterpart of the quantum-vacuum ``mock gravity'' of Spitzer and Gamow, whose cosmological
irrelevance was shown by Field \cite{Field1971}; in the \meq{} framework this exchange is the
\emph{local, always-present} mechanism.

\subsection{Orbital stability: the induced interaction must be Lorentz invariant}\label{sec:aberration}

Gravity propagates at $c$ through the substrate, consistent with GW170817
($|c_g - c|/c \lesssim 10^{-15}$ \cite{GW170817}). Any \emph{retarded flux} theory of a central
force would suffer Laplace-type aberration --- the force pointing to the retarded position of the
source, destabilizing orbits. Carlip \cite{Carlip2000} showed that in \emph{Lorentz-invariant}
field theories the aberration is cancelled to $\mathcal{O}(v/c)$ by velocity-dependent interaction
terms --- this is exactly how GR itself evades the Laplace argument. The \meq{} framework
therefore \textbf{requires} the Lorentz-invariant induced-interaction form of
Sec.~\ref{sec:induced}: the self-consistency equation must hold with the covariant structure
written there, and this is what guarantees orbital stability at the level verified by solar-system
dynamics and pulsar timing \cite{Nordtvedt1977}. This is a genuine, sharp constraint on the
microscopic model, and the induced-gravity form satisfies it by construction.

\subsection{Universality of free fall: coupling through \texorpdfstring{$T_{\mu\nu}$}{T-mu-nu}}

The framework couples gravity universally to total stress-energy $T_{\mu\nu}$ through the
self-consistency condition \eqref{eq:einstein} --- for ordinary matter (quasiparticle
excitations) and substrate background alike. The precision experiments then read as a
\emph{validation of this design choice at unprecedented level}:
\begin{itemize}
\item MICROSCOPE: $|\eta| < 2.1\times10^{-15}$ for Ti vs.\ Pt-Rh test masses \cite{Touboul2022};
\item neutron interferometry (Colella--Overhauser--Werner \cite{Colella1975}) and the
gravitational quantum states of ultracold neutrons (qBounce \cite{Jenke2014}) demonstrate the
coupling of individual fermions (the neutron) to Earth's gravity, with a phase/Bragg-type response
matching the equivalence principle;
\item the free fall of antihydrogen (ALPHA, 2023 \cite{ALPHA2023}) falls within GR.
\end{itemize}
The universal $T_{\mu\nu}$ coupling is also what protects the framework from the
composition-dependent ``fifth-force'' parameter region excluded by the E\"ot-Wash program
\cite{Lee2020}. Within our construction, spinor matter gravitates exactly as bosonic matter does.

\subsection{The gravitational quantum: a collective spin-2 mode}\label{sec:gravitonphoton}

In the \meq{} framework, the quantum of the gravitational field is a \textbf{collective mode of
the substrate} --- a quasiparticle with spin-2 quantum numbers, whose existence the framework
shares with all induced-gravity constructions. The spin-2 mode is not the photon; it is the
coherent density/current oscillation of the substrate, exactly as phonons are not photons
although both can scatter. This assignment is required by, and consistent with, four independent
lines of evidence:
\begin{enumerate}
\item \textbf{Spin.} The gravitational interaction couples to $T_{\mu\nu}$ with a spin-2
structure. The post-Newtonian phenomenology (perihelion precession, light deflection, Shapiro
delay, and the observed gravitational-wave orbital dynamics) is reproduced by a spin-2 field
\cite{Will2018} --- the substrate's collective spectrum must therefore contain this mode, and
the induced-gravity construction delivers it.
\item \textbf{Charge structure.} Electromagnetism has two signs of charge with both attraction
and repulsion and no universal coupling to neutral energy; gravity is universally attractive and
couples to all of $T_{\mu\nu}$. Distinct mediators are required --- precisely the collective
spin-0/½/1 (matter quasiparticles) and spin-2 (gravitational) modes of one substrate.
\item \textbf{Gravitational waves.} LIGO/Virgo events with electromagnetic counterparts
(GW170817) show that gravitational radiation propagates at $c$ and carries the tensor
polarization pattern of GR, while the associated photons arrive by a distinct channel with
different interactions --- consistent with distinct collective modes of the same substrate.
\item \textbf{Photon--photon scattering.} QED's Euler--Heisenberg nonlinearities --- confirmed
in Delbr\"uck scattering, light-by-light scattering at ATLAS, and the vacuum-birefringence
polarization signals of magnetars observed by IXPE \cite{ATLAS2019,PVLAS2016,Mignani2017} ---
saturate the observed photon--field interaction budget, bounding any additional substrate
contribution; this is the empirical basis of the $\zeta$-parameterization of Prediction 2
(Sec.~\ref{sec:p2}).
\end{enumerate}

\subsection{Cosmological transparency: bound on the inelastic \texorpdfstring{\meq{}--photon}{MEQ-photon} cross-section}

Photons propagate through the substrate without measurable energy loss, blurring, or anomalous
dispersion on cosmological baselines. Quantitatively, four independent observations constrain the
inelastic \meq{}--photon channel, written as a fractional energy loss $\xi$ per Hubble length:
\begin{enumerate}
\item \textbf{Time dilation of Type Ia supernova light curves} scales as $(1+z)$; the
spectral-aging rate follows $(1+z)^{0.97\pm0.10}$ \cite{Goldhaber2001} --- an energy-loss
mechanism produces no time dilation, so its contribution to the redshift budget is negligible.
\item \textbf{The Tolman surface-brightness test} gives $(1+z)^{-4}$ dimming, the signature of
expansion; HST cluster galaxy data verify it at high significance \cite{Lubin2001,Riess1998}.
\item \textbf{The CMB temperature scales as $T(z) = T_0(1+z)$} out to $z \approx 3$, as measured
from molecular excitation in high-$z$ absorbers --- the adiabatic cooling of an expanding photon
gas.
\item \textbf{The CMB spectrum is Planckian to $\sim$50 ppm (COBE/FIRAS)} \cite{Fixsen1996}, with
$|y| < 1.5\times10^{-5}$ on Compton distortions; any direction-randomizing or
energy-exchanging photon--medium interaction is excluded.
\end{enumerate}
The framework's conclusion is a design bound: \textbf{the \meq{}--photon inelastic cross-section
must be small enough that $\xi \lesssim 10^{-3}$ per Hubble length}, with the cosmological
redshift retaining its standard expansion interpretation and the accelerated expansion treated,
as in GR, as a property of the vacuum term $\Lambda$. This bound, together with the
elastic-channel constraints of Sec.~\ref{sec:gravitonphoton}, defines the transparency region of
parameter space in which Prediction 2 is formulated.

\subsection{Dark matter: the phase-dependent gravity conjecture}\label{sec:dm}

A distinctive conjecture of the \meq{} framework is \textbf{phase-dependent gravitation}:
gravitational strength tied to the \emph{dynamical phase} of the bound system --- systems in
formation (net inward substrate flux) exhibit a positive anomaly; systems in dissolution exhibit a
negative anomaly; mature systems exhibit none. The phenomenological target overlaps with MOND
\cite{McGaugh2016} (flat rotation curves without non-baryonic mass), but the mechanism is
different, and so are its signatures (Sec.~\ref{sec:p4}). Three classes of data constrain any
such modification:
\begin{enumerate}
\item \textbf{The Bullet Cluster} --- weak-lensing mass peaks displaced from the baryonic gas
\cite{Clowe2006} require the gravitating component to be collisionless relative to the gas,
which a pure field-phase effect tied to the baryon distribution struggles to produce.
\item \textbf{CMB acoustic peaks} --- the third-to-second peak ratio requires a
non-photon-coupled gravitating component at recombination \cite{Planck2020}.
\item \textbf{Big-bang nucleosynthesis and structure formation} --- constrain the
gravitating-to-baryonic density ratio at early times.
\end{enumerate}
A viable \meq{}-based alternative must therefore reproduce the \emph{full} gravitational
phenomenology of cold dark matter (lensing + CMB + structure formation), not merely rotation
curves. We do not claim this; we state it as the explicit burden of Prediction 4
(Sec.~\ref{sec:p4}) and note that the \meq{} signature --- gravitational strength depending on
the formation/dissolution phase of the bound system, strongest at early times --- is \emph{in
principle} distinguishable both from $\Lambda$CDM and from MOND by high-redshift rotation-curve
surveys with weak-lensing cross-checks, a program now becoming feasible with JWST.

\subsection{Short-range gravity and the substrate scale}\label{sec:shortrange}

Newton's inverse-square law is verified down to separations of $52\,\mu$m by the E\"ot-Wash
torsion-balance program \cite{Lee2020}, excluding gravity-strength Yukawa corrections with range
$\lambda > 39\,\mu$m (95\% C.L.). This bounds any substrate-scale modification of the
gravitational force to $\lambda \lesssim 40\,\mu$m or $|\alpha| \ll 1$ at larger $\lambda$.
Conversely, \textbf{the region below $52\,\mu$m is precisely where a substrate-scale effect must
live if it is to be observable at all.} The \meq{} framework therefore predicts, in its most
testable form, a deviation from the inverse-square law at separations of order the substrate
length $\ell_\ast < 40\,\mu$m, with strength set by the induced-gravity matching. This is a
genuine, falsifiable, laboratory-accessible prediction (Sec.~\ref{sec:p1}).

\subsection{Summary of consistency requirements}

\begin{table}[htbp]
\centering\small
\caption{Structural requirements of the \meq{} framework, their experimental basis, and status.}
\begin{tabular}{@{}lll@{}}
\toprule
Framework requirement & Experimental basis & Status \\
\midrule
Exact elasticity, detailed balance & No heating/drag anomalies & Satisfied by construction \\
Lorentz-invariant induced interaction & Orbital stability; GW170817 $c_g = c$ & Satisfied by construction \\
Universal coupling via $T_{\mu\nu}$ & MICROSCOPE $10^{-15}$; COW; qBounce; ALPHA & Validated at $10^{-15}$ \\
Spin-2 collective graviton & PPN phenomenology; tensor GW polarizations & Consistent \\
Cosmological transparency & SN~Ia dilation; Tolman; CMB $T(z)$, FIRAS & $\xi \lesssim 10^{-3}$/Hubble length \\
Substrate length $\ell_\ast$ & E\"ot-Wash: $\lambda<39\,\mu$m excluded at $|\alpha|\!\sim\!1$ & $\ell_\ast < 40\um$ $\to$ \textbf{testable} (P1) \\
Photon--field elastic channel & IXPE; PVLAS; light-by-light & $|\zeta| \lesssim 10^{-3}$--$10^{-1}$ $\to$ \textbf{testable} (P2) \\
Phase-dependent gravitation & Bullet Cluster; CMB peaks; BBN & Constrained $\to$ \textbf{testable} (P4) \\
\bottomrule
\end{tabular}
\end{table}

\section{Novel Predictions and the Experimental Program}\label{sec:pred}

\subsection{Summary of the framework's novel content}

The \meq{} framework, in the form established by Secs.~\ref{sec:concept}--\ref{sec:exp},
contributes: (i) the operational derivation of emergent, light-speed-defined spacetime
(Sec.~\ref{sec:concept}); (ii) the induced-gravity self-consistency structure with
higher-curvature corrections suppressed by the substrate scale (Sec.~\ref{sec:induced});
(iii) the topological program for fermionic quasiparticles (Sec.~\ref{sec:fermions});
(iv) emergent Lorentz symmetry with tiny violations near $E_\ast$ (Sec.~\ref{sec:kinetic});
and (v) a kinetic (exchange-based) interpretation of the gravitational interaction that yields a
\emph{mass non-stationarity} prediction no field theory of gravity makes
(Sec.~\ref{sec:p3} below). We now state the four falsifiable predictions.

\subsection{Prediction 1: sub-40-\texorpdfstring{$\mu$}{mu}m deviation from the inverse-square law}\label{sec:p1}

The framework predicts a Yukawa-type deviation from Newtonian gravity,
\begin{equation}
V(r) \;=\; -\frac{G m_1 m_2}{r}\Bigl[\,1 + \alpha\, e^{-r/\lambda}\Bigr],
\label{eq:yukawa}
\end{equation}
with $\lambda \approx \ell_\ast$ (the substrate length) and $\alpha$ of order unity (naturalness
of the induced-gravity matching). Current data \cite{Lee2020} exclude this for $\lambda >
39\,\mu$m at $|\alpha| \sim 1$.

\textbf{Quantitative estimate (details in Appendix~\ref{app:quant}).} For two parallel tungsten
plates at gap $d$ the Yukawa pressure is
\begin{equation}
P_Y(d) \;=\; 2\pi\, G\, \alpha\, \lambda^{2} \rho_1 \rho_2\, e^{-d/\lambda},
\label{eq:pyuk}
\end{equation}
numerically $P_Y \approx 3.3\times10^{-11}\,$Pa at the current exclusion boundary
($\alpha=1$, $\lambda=39\,\mu$m, $d=52\,\mu$m) --- i.e., the E\"ot-Wash 2020 experiment
operates at a pressure sensitivity of order $10^{-11}\,$Pa. A next-generation experiment
($A = 100\,$cm$^2$ of W--W working area, gap $40\,\mu$m) with force noise
$F_{\min} = 10^{-16}$--$10^{-17}\,$N reaches
$|\alpha|_{\min} \approx 10^{-3}$--$10^{-4}$ at $\lambda = 20\,\mu$m
(Fig.~\ref{fig:yukawa}). The Casimir background at these gaps
($\sim\!10^{-7}\,$Pa at $10\,\mu$m) exceeds the signal but has a distinct $d^{-4}$ versus
$e^{-d/\lambda}$ separation dependence, allowing discrimination. \textbf{Prediction:} a
deviation appears as experiments push below $\sim40\,\mu$m toward $\ell_\ast$; otherwise the
bound pushes $\ell_\ast$ toward its Planck-scale limit and the substrate becomes empirically
indistinguishable from quantum-gravity EFT.

\begin{figure}[htbp]
\centering
\includegraphics[width=0.85\textwidth]{fig1_yukawa_reach.png}
\caption{Discovery reach in the $(\alpha,\lambda)$ plane for Prediction 1. Red: approximate
E\"ot-Wash 2020 exclusion; green/blue dotted: reach of a next-generation experiment with
$A=100\,$cm$^2$ W--W at gap $40\,\mu$m for two force-noise floors. The region between the curves
is the viable discovery window of the \meq{} substrate length $\ell_\ast$.}
\label{fig:yukawa}
\end{figure}

\subsection{Prediction 2: anomalous photon--field interactions below the QED floor}\label{sec:p2}

The substrate necessarily contributes a forward-scattering amplitude to photon propagation
through strong ambient fields, of the same parametric form as the Euler--Heisenberg correction
but with a different coefficient $\zeta$ (and no electron-mass Compton-scale suppression if
$\ell_\ast \neq \lambda_C$). Writing the anomalous vacuum-birefringence coefficient as
$\zeta \times (\text{QED value})$, current bounds from PVLAS, OSQAR, gamma-ray astrophysics, and
magnetar polarimetry \cite{PVLAS2016,Mignani2017} require $|\zeta| \lesssim 10^{-3}$--$10^{-1}$
depending on channel.

\textbf{Quantitative estimate (Appendix~\ref{app:quant}).} The Euler--Heisenberg Cotton--Mouton
asymmetry is $\Delta n_{\rm QED} = \frac{\alpha_{\rm em}}{45\pi}(B/B_c)^2$ with $B_c =
4.414\times10^{9}\,$T; at $B = 10\,$T this is $\Delta n = 2.6\times10^{-22}$. The accumulated
ellipticity in a cavity is $\psi = N_{\rm eff}\,\pi\,\Delta n\, L/\lambda_{\rm light}$: for the
operated PVLAS parameters ($B=2.5\,$T, $L=1\,$m, $N_{\rm eff}=2.5\times10^5$) the QED signal is
$\psi_{\rm QED} \approx 1.2\times10^{-11}\,$rad --- the QED floor is \emph{itself} at the edge
of detectability, which is why any $\zeta \gtrsim 0.1$ would already have been seen. A
next-generation setup (13--14\,T magnet, 10--25\,m cavity, $N_{\rm eff} \sim 10^{5}$) delivers
$\psi_{\rm QED} \sim 10^{-9}\,$rad, so a readout noise of $10^{-11}\,$rad probes
$\zeta \sim 10^{-2}$--$10^{-3}$ (Fig.~\ref{fig:biref}). \textbf{Prediction:} a residual,
polarization-dependent phase shift in multi-pass optical cavity experiments embedded in strong
($\geq$10\,T) laboratory fields, at the level set by $\zeta \neq 0$, observable \emph{below} the
QED floor once the QED signal itself is calibrated. A null result down to
$\zeta \sim 10^{-6}$ would force the \meq{}--photon coupling into the same structure as ordinary
QED loops.

\begin{figure}[htbp]
\centering
\includegraphics[width=0.95\textwidth]{fig2_birefringence.png}
\caption{Prediction 2. Left: Euler--Heisenberg birefringence $\Delta n_{\rm QED}(B)$. Right:
cavity ellipticity versus effective passes for PVLAS-class and next-generation (BMV-class)
setups; the horizontal line marks a representative readout noise floor.}
\label{fig:biref}
\end{figure}

\subsection{Prediction 3: mass non-stationarity of gravitating configurations}\label{sec:p3}

Because the gravitational field is maintained by the balance of substrate absorption and
emission, an exact prediction of the kinetic picture is that the \emph{inertial mass of a body
is stationary only in the detailed-balance limit}. The momentum flux that sustains the weight
$F = mg$ of a body carries energy flux $\sim\! Fc = mgc$; if a fraction $\eta$ of that anisotropic
power is stored as mass rather than re-emitted, then
\begin{equation}
\frac{1}{m}\frac{dm}{dt} \;=\; \eta\,\frac{g}{c},
\label{eq:etadrift}
\end{equation}
where $g$ is the local gravitational field. At Earth's surface $g/c = 3.3\times10^{-8}\,$s$^{-1}$
($\approx 1.03\,$yr$^{-1}$ for $\eta=1$), so existing data already constrain $\eta$ severely.

\textbf{Quantitative estimate (Appendix~\ref{app:quant}).} Converting existing limits on
anomalous mass drift:
\begin{itemize}
\item \emph{Binary pulsars.} The Hulse--Taylor pulsar's orbital decay agrees with GR to
$|\dot P_b^{\rm anom}| \lesssim 1.6\times10^{-14}\,$s/s; with $\dot P/P = -2\,\dot m/m$ for
adiabatic isotropic mass change, and local field $g = 50\,$m/s$^2$ at the pulsar,
$\eta \lesssim 4.8\times10^{-8}$. The double pulsar J0737$-$3039 (conservatively
$10^{-3}\times$GR in $\dot P_b$, $g = 214\,$m/s$^2$ \cite{Kramer2021}) gives
$\eta \lesssim 9\times10^{-10}$.
\item \emph{Laboratory mass metrology.} Composition-\emph{differential} stability of standards
($\sim\!5\times10^{-9}$/yr) implies $\eta_{\rm comp} \lesssim 5\times10^{-9}$.
\item \emph{Atomic-clock comparisons.} Limits on variation of mass ratios
($\lesssim 10^{-16}$/yr) imply $\eta_{\rm comp} \lesssim 10^{-16}$ \emph{if} the drift is
carried by composition-relevant masses --- the strongest existing constraint on the
composition-dependent part.
\end{itemize}
\textbf{The untested content is a common-mode, composition-\emph{independent}} drift, which
pulsar timing bounds only at $\eta \lesssim 10^{-9}$, plus any \emph{potential-scaled} drift
$\dot m/m = \Gamma\,\Phi/c^{2}$, whose distinctive signature is an \emph{annual modulation} of
mass ratios (amplitude $\Gamma \times 1.7\times10^{-3}\,$s) synchronized with the Earth's orbital
phase ($\Delta\Phi/c^2 = \pm3.3\times10^{-10}$ from eccentricity $e = 0.0167$).
\textbf{Falsifiable form:} a differential mass-drift experiment monitoring two
composition-distinct test masses (e.g., Si vs.\ Pt-Ir) with $10^{-12}$--$10^{-13}$/yr ratio
precision and Fourier analysis at $f = 1/$yr against the known phase of Earth's orbit probes
$\eta_{\rm diff} \sim 10^{-8}$--$10^{-11}$ --- unexplored territory. A positive signal would be a
unique signature impossible for any field theory of gravity (in which inertial mass is exactly
conserved), while a null result bounds the exchange asymmetry. We emphasize that this prediction
follows only from the \emph{kinetic} (exchange-based) reading of the framework; the pure
induced-gravity EFT reading predicts no such drift. The experiment therefore \emph{discriminates
between the two readings} --- a logically clean and, to our knowledge, novel test.

\begin{figure}[htbp]
\centering
\includegraphics[width=0.9\textwidth]{fig3_eta_landscape.png}
\caption{Prediction 3: landscape of bounds on the exchange-imbalance parameter $\eta$ from
binary pulsars, laboratory mass metrology, and clock comparisons, together with the projected
reach of the proposed differential mass-comparison experiment.}
\label{fig:eta}
\end{figure}

\subsection{Prediction 4: phase-dependent galactic gravitation}\label{sec:p4}

If the phase-dependent mechanism of Sec.~\ref{sec:dm} has any macroscopic content, the most
distinctive signature is gravitational strength tied to the \emph{dynamical phase} of the bound
system: systems in formation (net inward flux) exhibit a positive anomaly; systems in dissolution
exhibit a negative anomaly; mature systems exhibit none. This differs from MOND (which depends
on acceleration scale $a_0$) and from $\Lambda$CDM (which depends on halo history).
\textbf{Test:} joint rotation-curve + weak-lensing surveys of galaxies as a function of redshift
and star-formation phase; the \meq{} signature is a correlation of the ``missing mass'' ratio
with specific star-formation rate and environment at fixed stellar mass, which $\Lambda$CDM
predicts to be weak and MOND predicts to be absent. Existing data (e.g., the SPARC sample at low
$z$) already provide partial leverage \cite{McGaugh2016}; JWST and ELT high-$z$ surveys complete
it. We stress that this prediction survives only in the region of parameter space where the
Bullet Cluster and CMB constraints are satisfied by the baryonic-plus-phase content --- a
demanding requirement that may yet exclude the mechanism entirely; we state it as a
\emph{target for falsification}, not as an achieved result.

\section{Open Problems and Limitations}\label{sec:open}

We state frankly the points at which the program is incomplete:
\begin{enumerate}
\item \textbf{Fermion construction.} No 3+1-dimensional topological construction of the full
Standard Model fermion spectrum from a single bosonic substrate exists. The Skyrme mechanism
demonstrates existence-in-principle for single fermions, not a spectrum with the observed
chiral/gauge structure.
\item \textbf{Gauge group emergence.} The emergence of the full Standard Model gauge group
$\mathrm{U}(1)\times\mathrm{SU}(2)\times\mathrm{SU}(3)$ from the substrate is not constructed
here; the established route in substrate models is via the topology of the vacuum manifold
(Volovik's program \cite{Volovik2003}), and this is the path the program must follow.
\item \textbf{Quantitative cosmology.} Inflation, baryogenesis, and the CMB power spectrum have
no treatment in the framework. If phase-dependent gravity (Sec.~\ref{sec:p4}) fails the
Bullet-Cluster class of tests, the framework must revert to $\Lambda$CDM cosmology wholesale and
its distinctive content shrinks to laboratory-scale effects
(Secs.~\ref{sec:p1}--\ref{sec:p3}).
\item \textbf{Parameter naturalness.} The framework survives only with
$\ell_\ast < 40\,\mu$m, $\xi$ (inelastic photon loss) $< 10^{-3}$ per Hubble length, and
$\zeta < 10^{-3}$. Each of these can be viewed as a fine-tuning (a criticism) or as a set of
falsifiable restrictions (our reading), but honesty requires both readings to be stated.
\end{enumerate}

The honest overall assessment is: \textbf{the \meq{} framework, in the form presented here, is a
member of the family of induced/emergent gravity models, distinguished by its kinetic (exchange)
interpretation and its consequent non-stationarity prediction (Sec.~\ref{sec:p3}); it is not
presently a competitor to established theory in explanatory scope, and its value is contingent
on the outcomes of the tests in Sec.~\ref{sec:pred}.}

\section{Conclusion}\label{sec:conclusion}

We have formulated the Minimum Energy Quantum framework --- a single-species, light-speed,
elastic-exchange substrate from which the operational content of spacetime, the relativity
postulates, an induced gravitational interaction with spin-2 collective quanta, and a program
for the emergence of matter are to be derived --- and established its consistency with the
complete corpus of modern experimental constraints, from MICROSCOPE and E\"ot-Wash through
GW170817 to the CMB. The framework's four novel elements --- the operational derivation of
relativistic spacetime, the induced-gravity self-consistency structure, the topological program
for fermions, and emergent Lorentz symmetry --- yield four falsifiable predictions with
quantitative reach estimates: sub-40-$\mu$m inverse-square-law deviations (with the
$(\alpha,\lambda)$ discovery reach of next-generation torsion-balance experiments computed),
sub-QED-floor photon--field interactions (with cavity ellipticities computed for current and
next-generation setups), composition-dependent mass drift with annual-potential tagging (with
bounds from binary pulsars, mass metrology, and clock comparisons), and phase-dependent galactic
gravitation --- each with identified experiments capable of deciding them within the current
decade. Whatever the outcomes, the framework will then have done what a scientific hypothesis is
supposed to do: expose itself to decisive evidence.

\appendix

\section{Quantitative Estimates}\label{app:quant}

This appendix collects the numerical estimates behind Predictions 1--3. All values were computed
from the standard constants ($G = 6.674\times10^{-11}\,$SI, $c = 2.998\times10^{8}\,$m/s,
$\hbar$, $\alpha_{\rm em}$, $B_c = m_e^2c^2/e\hbar = 4.414\times10^{9}\,$T); the computation
script accompanies the paper as supplementary material.

\subsection*{A.1\ \ Yukawa pressure and discovery reach (Prediction 1)}

For the parallel-plate geometry, Eq.~\eqref{eq:pyuk} gives, for tungsten
($\rho = 1.93\times10^{4}\,$kg/m$^3$), $\alpha = 1$:

\begin{table}[htbp]
\centering\small
\caption{Yukawa pressure between parallel W plates, $\alpha=1$ (from Eq.~\eqref{eq:pyuk}).}
\begin{tabular}{@{}rrrr@{}}
\toprule
$\lambda$ [$\mu$m] & $P_Y(d=\lambda)$ [Pa] & $P_Y(d=2\lambda)$ [Pa] & $P_Y(d=52\,\mu$m) [Pa]\\
\midrule
5  & $1.4\times10^{-12}$ & $5.3\times10^{-13}$ & $1.2\times10^{-16}$ \\
10 & $5.7\times10^{-12}$ & $2.1\times10^{-12}$ & $8.6\times10^{-14}$ \\
20 & $2.3\times10^{-11}$ & $8.5\times10^{-12}$ & $4.6\times10^{-12}$ \\
30 & $5.2\times10^{-11}$ & $1.9\times10^{-11}$ & $2.5\times10^{-11}$ \\
40 & $9.2\times10^{-11}$ & $3.4\times10^{-11}$ & $6.8\times10^{-11}$ \\
60 & $2.1\times10^{-10}$ & $7.6\times10^{-11}$ & $2.4\times10^{-10}$ \\
\bottomrule
\end{tabular}
\end{table}

\noindent
The E\"ot-Wash 2020 limit ($\alpha = 1$ excluded for $\lambda > 39\,\mu$m at working gap
$d = 52\,\mu$m) corresponds to a detected pressure $\sim\!3\times10^{-11}\,$Pa. With $A =
100\,$cm$^2$ of active area at gap $40\,\mu$m, force noise floors of $10^{-15}$/$10^{-16}$/
$10^{-17}\,$N correspond to $|\alpha|_{\min}(\lambda = 20\,\mu{\rm m}) \approx 1.2\times
10^{-2}/10^{-3}/10^{-4}$. The Casimir pressure $\pi^2\hbar c/240d^4 = 1.3\times10^{-7}\,$Pa at
$d = 10\,\mu$m exceeds a gravity-strength Yukawa signal ($5.7\times10^{-12}\,$Pa) by
$\sim\!2\times10^{4}$ but scales as $d^{-4}$, enabling separation by gap scans.

\subsection*{A.2\ \ Vacuum birefringence and the $\zeta$-reach (Prediction 2)}

$\Delta n_{\rm QED}(B) = (\alpha_{\rm em}/45\pi)(B/B_c)^2 = 2.65\times10^{-24}\,B^2[\mathrm{T}^2]$;
$\psi = N_{\rm eff}\,\pi\Delta n L/\lambda_{\rm light}$ with $\lambda_{\rm light} = 1064\,$nm:

\begin{table}[htbp]
\centering\small
\caption{Cavity ellipticities of the QED signal and the $\zeta$ reach.}
\begin{tabular}{@{}lrrrrc@{}}
\toprule
Setup & $B$ [T] & $L$ [m] & $N_{\rm eff}$ & $\psi_{\rm QED}$ [rad] & $\zeta_{\min}$ ($10^{-11}\,$rad noise)\\
\midrule
PVLAS (operated)  & 2.5 & 1.0 & $2.5\times10^{5}$ & $1.2\times10^{-11}$ & $\sim\!1$ \\
PVLAS (upgraded)   & 2.5 & 3.2 & $4.5\times10^{5}$ & $7.0\times10^{-11}$ & $1.4\times10^{-1}$ \\
OSQAR-class        & 9.0 & 14.3 & $10^{2}$ & $9.1\times10^{-13}$ & $1.1\times10^{1}$ \\
BMV-class          & 13  & 25 & $6.4\times10^{4}$ & $2.1\times10^{-9}$ & $4.7\times10^{-3}$ \\
\bottomrule
\end{tabular}
\end{table}

\subsection*{A.3\ \ Mass non-stationarity bounds (Prediction 3)}

Equation~\eqref{eq:etadrift} with existing limits: Hulse--Taylor
($|\dot P_b^{\rm anom}|<1.6\times10^{-14}\,$s/s, $g = 50.3\,$m/s$^2$)
$\Rightarrow \eta < 4.8\times10^{-8}$; double pulsar J0737$-$3039 (conservative
$10^{-3}\times$GR, $g = 214\,$m/s$^2$) $\Rightarrow \eta < 8.7\times10^{-10}$; laboratory
composition-differential mass stability ($5\times10^{-9}$/yr) $\Rightarrow
\eta_{\rm comp} < 4.8\times10^{-9}$; clock-ratio stability ($10^{-16}$/yr) $\Rightarrow
\eta_{\rm comp} < 10^{-16}$. Annual modulation of a potential-scaled drift
$\dot m/m = \Gamma\,\Phi/c^2$ has amplitude $\Gamma\times1.7\times10^{-3}\,$s at $f = 1/$yr;
differential monitoring at $10^{-12}$ probes $\Gamma < 6\times10^{-10}\,$s$^{-1}$. The proposed
differential experiment (Sec.~\ref{sec:p3}) probes $\eta_{\rm diff} \sim 10^{-8}$--$10^{-11}$,
between the pulsar and clock bounds --- currently unexplored.

\begin{thebibliography}{50}\small

\bibitem{Rovelli2015} Rovelli, C. \& Vidotto, F. \emph{Covariant Loop Quantum Gravity: An
Elementary Introduction to Quantum Gravity and Spinfoam Theory} (Cambridge University Press,
2015).

\bibitem{Polchinski1998} Polchinski, J. \emph{String Theory} (Cambridge University Press, 1998);
Maldacena, J. The large-$N$ limit of superconformal field theories and supergravity.
\emph{Adv. Theor. Math. Phys.} \textbf{2}, 231 (1998).

\bibitem{Reuter1998} Reuter, M. Nonperturbative evolution equation for quantum gravity.
\emph{Phys. Rev. D} \textbf{57}, 971 (1998); Eichhorn, A. et al. Status reviews in
\emph{Phys. Rep.} (2022).

\bibitem{Loll2020} Loll, R. Quantum gravity from causal dynamical triangulations: a review.
\emph{Class. Quantum Grav.} \textbf{37}, 243002 (2020).

\bibitem{Batista2025} Alves Batista, R. et al. White paper and roadmap for quantum gravity
phenomenology in the multi-messenger era. \emph{Class. Quantum Grav.} \textbf{42}, 032001 (2025).

\bibitem{Touboul2022} Touboul, P. et al. (MICROSCOPE Collaboration). Result of the MICROSCOPE
weak equivalence principle test. \emph{Class. Quantum Grav.} \textbf{39}, 204009 (2022).

\bibitem{Williams2004} Williams, J. G., Turyshev, S. G. \& Boggs, D. H. Progress in lunar laser
ranging tests of relativistic gravity. \emph{Phys. Rev. Lett.} \textbf{93}, 261101 (2004);
Nordtvedt, K. Testing relativity with laser ranging to the Moon. \emph{Phys. Rev.}
\textbf{170}, 1186 (1968).

\bibitem{GW170817} Abbott, B. P. et al. Gravitational waves and gamma-rays from a binary neutron
star merger: GW170817 and GRB 170817A. \emph{Astrophys. J. Lett.} \textbf{848}, L13 (2017).

\bibitem{Abdo2009} Abdo, A. A. et al. (Fermi GBM/LAT). A limit on the variation of speed of
light arising from Fermi observations of GRB~090510. \emph{Science} \textbf{323}, 1688 (2009).

\bibitem{Vasileiou2013} Vasileiou, V. et al. Constraints on Lorentz invariance violation from
Fermi-Large Area Telescope observations of gamma-ray bursts. \emph{Phys. Rev. D} \textbf{87},
122001 (2013).

\bibitem{Jacobson2003} Jacobson, T., Liberati, S. \& Mattingly, D. Lorentz invariance and
ultra-high energy cosmic rays. \emph{Phys. Rev. D} \textbf{67}, 124011 (2003).

\bibitem{Myers2005} Myers, R. C. \& Pospelov, M. Experimental challenges for quantum gravity.
\emph{Phys. Rev. Lett.} \textbf{90}, 211601 (2003).

\bibitem{Bolokhov2007} Bolokhov, P. A. \& Pospelov, M. Classification of dimension 5 Lorentz
violating interactions in the standard model. \emph{Class. Quantum Grav.} \textbf{24}, 4991
(2007).

\bibitem{LeSage2002} Edwards, M. R. (ed.). \emph{Pushing Gravity: New Perspectives on Le Sage's
Theory of Gravitation} (Apeiron, 2002).

\bibitem{Poincare1908} Poincar\'e, H. La dynamique de l'\'electron. \emph{Rev. G\'en\'erale des
Sci.} \textbf{19}, 386 (1908); Maxwell, J. C. Atom. \emph{Encyclop\ae dia Britannica} (1875).

\bibitem{Sakharov1968} Sakharov, A. D. Vacuum quantum fluctuations in curved space and the
theory of gravitation. \emph{Sov. Phys. Dokl.} \textbf{12}, 1040 (1968).

\bibitem{Donoghue1994} Donoghue, J. F. General relativity as an effective field theory: the
leading quantum corrections. \emph{Phys. Rev. D} \textbf{50}, 3874 (1994).

\bibitem{Jacobson1995} Jacobson, T. Thermodynamics of spacetime: the Einstein equation of state.
\emph{Phys. Rev. Lett.} \textbf{75}, 1260 (1995).

\bibitem{Verlinde2011} Verlinde, E. On the origin of gravity and the laws of Newton.
\emph{JHEP} \textbf{1104}, 029 (2011); Padmanabhan, T. Thermodynamical aspects of gravity.
\emph{Rep. Prog. Phys.} \textbf{73}, 046901 (2010).

\bibitem{Volovik2003} Volovik, G. E. \emph{The Universe in a Helium Droplet} (Oxford University
Press, 2003).

\bibitem{Barcelo2005} Barcel\'o, C., Liberati, S. \& Visser, M. Analogue gravity.
\emph{Living Rev. Rel.} \textbf{8}, 12 (2005).

\bibitem{Sinha1976} Sinha, K. P., Sivaram, C. \& Sudarshan, E. C. G. A superfluid vacuum theory of
matter and gravity. \emph{Found. Phys.} \textbf{6}, 65 (1976).

\bibitem{Zloshchastiev2011} Zloshchastiev, K. G. Spontaneous symmetry breaking in a logarithmic
nonlinear quantum theory. \emph{Acta Phys. Polon. B} \textbf{42}, 261 (2011).

\bibitem{Mansouri1977} Mansouri, R. \& Sexl, R. U. A test theory of special relativity.
\emph{Gen. Rel. Grav.} \textbf{8}, 497 (1977).

\bibitem{Amelino2002} Amelino-Camelia, G. Doubly-special relativity. \emph{Nature} \textbf{418},
34 (2002).

\bibitem{Skyrme1962} Skyrme, T. H. R. A unified field theory of mesons and baryons.
\emph{Nucl. Phys.} \textbf{31}, 556 (1962).

\bibitem{Goldstone1981} Goldstone, J. \& Wilczek, F. Fractional quantum numbers on solitons.
\emph{Phys. Rev. Lett.} \textbf{47}, 986 (1981); Wilczek, F. \& Zee, A. Linking numbers, spin,
and statistics of solitons. \emph{Phys. Rev. Lett.} \textbf{51}, 2250 (1983).

\bibitem{Wilczek1982} Wilczek, F. Magnetic flux, angular momentum, and statistics.
\emph{Phys. Rev. Lett.} \textbf{48}, 1144 (1982).

\bibitem{Field1971} Field, G. B. The time relaxation of a resonance-line photon gas.
\emph{Astrophys. J.} \textbf{167}, 231 (1971).

\bibitem{Carlip2000} Carlip, S. Aberration and the speed of gravity. \emph{Phys. Lett. A}
\textbf{267}, 81 (2000).

\bibitem{Nordtvedt1977} Nordtvedt, K. \& Will, C. M. Testing relativity with pulsar timing.
\emph{Astrophys. J.} \textbf{212}, 475 (1977).

\bibitem{Colella1975} Colella, R., Overhauser, A. W. \& Werner, S. A. Observation of
gravitationally induced quantum interference. \emph{Phys. Rev. Lett.} \textbf{34}, 1472 (1975).

\bibitem{Jenke2014} Jenke, T. et al. (qBounce). Measurement of gravitational states of ultracold
neutrons in the Earth's gravitational field. \emph{Phys. Rev. Lett.} \textbf{112}, 151105 (2014);
Cronenberg, G. et al. Gravity resonance spectroscopy constrains dark energy and dark matter
scenarios. \emph{Nature Phys.} \textbf{14}, 1022 (2018).

\bibitem{ALPHA2023} Anderson, E. K. et al. (ALPHA Collaboration). Observation of the effect of
gravity on the motion of antimatter. \emph{Nature} \textbf{621}, 716 (2023).

\bibitem{Lee2020} Lee, J. G. et al. (E\"ot-Wash Group). New test of the gravitational $1/r^2$ law
at separations down to 52~$\mu$m. \emph{Phys. Rev. Lett.} \textbf{124}, 101101 (2020);
Adelberger, E., Gundlach, J., Heckel, B., Hoedl, S. \& Schlamminger, S. Torsion balance
experiments: a low-energy frontier of fundamental physics. \emph{Prog. Part. Nucl. Phys.}
\textbf{62}, 102 (2009).

\bibitem{Will2018} Will, C. M. The confrontation between general relativity and experiment.
\emph{Living Rev. Rel.} \textbf{17}, 4 (2014); \emph{Theory and Experiment in Gravitational
Physics} (Cambridge University Press, 2018).

\bibitem{ATLAS2019} ATLAS Collaboration. Observation of light-by-light scattering in
ultraperipheral Pb+Pb collisions with the ATLAS detector. \emph{Phys. Rev. Lett.} \textbf{123},
051801 (2019).

\bibitem{PVLAS2016} Della Valle, F. et al. (PVLAS). The PVLAS experiment: measuring vacuum
magnetic birefringence and dichroism with an upgraded magnetically induced optical cavity.
\emph{Eur. Phys. J. C} \textbf{76}, 24 (2016).

\bibitem{Mignani2017} Mignani, R. P. et al. Evidence for vacuum birefringence from the first
optical-polarimetry measurement of the isolated neutron star RX~J1856.5$-$3754. \emph{Mon. Not.
R. Astron. Soc.} \textbf{465}, 492 (2017).

\bibitem{Goldhaber2001} Goldhaber, G. et al. Timescale stretch parameterization of Type~Ia
supernova $B$-band light curves. \emph{Astrophys. J.} \textbf{558}, 359 (2001).

\bibitem{Lubin2001} Lubin, L. M. \& Sandage, A. The Tolman surface brightness test for the
reality of the expansion. IV. A measurement of the Tolman signal and the luminosity evolution.
\emph{Astron. J.} \textbf{122}, 1084 (2001).

\bibitem{Riess1998} Riess, A. G. et al. Observational evidence from supernovae for an
accelerating universe and a cosmological constant. \emph{Astron. J.} \textbf{116}, 1009 (1998);
Perlmutter, S. et al. Measurements of $\Omega$ and $\Lambda$ from 42 high-redshift supernovae.
\emph{Astrophys. J.} \textbf{517}, 565 (1999).

\bibitem{Fixsen1996} Fixsen, D. J. et al. (COBE/FIRAS). The cosmic microwave background spectrum
from the full COBE FIRAS data set. \emph{Astrophys. J.} \textbf{473}, 576 (1996).

\bibitem{McGaugh2016} McGaugh, S. S., Lelli, F. \& Schombert, J. I. Radial acceleration relation
in rotationally supported galaxies. \emph{Phys. Rev. Lett.} \textbf{117}, 201101 (2016).

\bibitem{Clowe2006} Clowe, D. et al. A direct empirical proof of the existence of dark matter.
\emph{Astrophys. J. Lett.} \textbf{648}, L109 (2006).

\bibitem{Planck2020} Planck Collaboration. Planck 2018 results. VI. Cosmological parameters.
\emph{Astron. Astrophys.} \textbf{641}, A6 (2020).

\bibitem{Kramer2021} Kramer, M. et al. Tests of relativistic gravity with radio pulsars.
\emph{Ann. Rev. Nucl. Part. Sci.} \textbf{73}, 1 (2023).

\end{thebibliography}

\end{document}
